\section{Architecture}

% Slide 6: Proposed Methodology
\begin{frame}{Proposed Methodology}
  We introduce the \textbf{Synapse Dynamic Mesh} framework:
  \begin{enumerate}
    \item \textbf{Asynchronous Spike Routing:} Messages are sent only when neuron potentials exceed threshold $\theta_i$.
    \item \textbf{Dynamic Plasticity Engine:} Synaptic weights are updated locally using an optimized Hebbian learning model.
    \item \textbf{Adaptive Power Scaling:} Sleep states are automatically triggered for idle neural core clusters.
    \item \textbf{Decentralized Consensus:} Multi-agent coordination allows isolated sub-networks to pool learned insights without central coordination.
  \end{enumerate}
\end{frame}

% Slide 7: System Workflow
\begin{frame}{System Workflow Diagram}
  \begin{center}
    \begin{tikzpicture}[node distance=1.2cm, auto, thick, scale=0.85, every node/.style={scale=0.85}]
      % Nodes
      \node [draw, fill=blue!10, rectangle, rounded corners, minimum width=2cm, minimum height=0.8cm, align=center] (input) {Raw Input \\ Data Feed};
      \node [draw, fill=teal!10, rectangle, rounded corners, minimum width=2.2cm, minimum height=0.8cm, right=of input, align=center] (extract) {Feature \\ Extractor};
      \node [draw, fill=orange!10, rectangle, rounded corners, minimum width=2.4cm, minimum height=0.8cm, right=of extract, align=center] (process) {Cognitive Core \\ (Synapse Engine)};
      \node [draw, fill=green!10, rectangle, rounded corners, minimum width=2cm, minimum height=0.8cm, right=of process, align=center] (output) {Decision \\ Output};
      
      % Arrows
      \draw [->, >=stealth, thick] (input) -- (extract);
      \draw [->, >=stealth, thick] (extract) -- (process);
      \draw [->, >=stealth, thick] (process) -- (output);
      
      % Descriptions
      \node [below=0.6cm of extract, text width=3.5cm, align=center, font=\scriptsize\itshape] (desc1) {Normalizes signals to spike-train streams};
      \node [below=0.6cm of process, text width=3.5cm, align=center, font=\scriptsize\itshape] (desc2) {Updates dynamic weights via STDP rule};
      
      \draw [dashed, ->] (extract) -- (desc1);
      \draw [dashed, ->] (process) -- (desc2);
    \end{tikzpicture}
  \end{center}
  \small
  The flow shows signals entering from sensors, being transformed to temporal spike trains, processed by our plastic SNN, and translated to actions.
\end{frame}

% Slide 8: Multi-Agent Coordination
\begin{frame}{Multi-Agent Coordination}
  \begin{columns}[T]
    \begin{column}{0.48\textwidth}
      \textbf{Core Concepts:}
      \vspace{0.2cm}
      
      Each agent is an independent Synapse instance deployed on local hardware. 
      
      Agents communicate via a lightweight, asynchronous mesh protocol (SynapseMesh) to coordinate routing decisions.
    \end{column}
    \begin{column}{0.48\textwidth}
      \textbf{Key Characteristics:}
      \begin{itemize}
        \item \textbf{Decentralized:} No single point of failure.
        \item \textbf{Low Overhead:} Message sizes are restricted to 64 bytes.
        \item \textbf{Resilient:} Robust to up to 30\% agent offline rates.
        \item \textbf{Scalable:} Linear complexity scaling.
      \end{itemize}
    \end{column}
  \end{columns}
\end{frame}

% Section 3: Performance
